%%============================================================================
%% A report. Created by report-new.sh. Do not copy this file by hand.
%% Compile from the repository root:  report-build.sh lib/reports/src/<this file>
%%
%% Two rules the build enforces and you cannot argue with.
%%   Section 0, when this repository asks for one, lives in
%%   lib/reports/reflections/<same stem>.tex and the author writes it.
%%   The PDF is five pages, references and appendix included.
%%============================================================================
\documentclass{labreport}
\lrnoreflection   % lib/reports/src/config.yml says reflection: false
\usepackage{amssymb}

\lrdate{2026-10-16}
\lrslug{halil-learning-oracle-definition}
\lrtitle{A learning oracle for procedures that train and query models}
\lrauthor{Halil İbrahim Kanpak}
\lrstatus{draft}          % draft | reviewed | final

\lrdefine{MnistCifar}{215}{lib/reports/src/data/2026-10-16-halil-swept-papers-and-threat-models-facts.csv}
\lrdefine{SweptPapers}{349}{lib/reports/src/data/2026-10-16-halil-swept-papers-and-threat-models-facts.csv}

\newcommand{\Z}{\mathcal{Z}}
\newcommand{\Mc}{\mathcal{M}}
\newcommand{\Hc}{\mathcal{H}}
\newcommand{\HH}{\mathbb{H}}
\newcommand{\Lf}{\mathsf{L}}
\newcommand{\ev}{\mathsf{ev}}
\newcommand{\Or}{\mathcal{O}}
\newcommand{\Ob}{\widehat{\mathcal{O}}}
\newcommand{\Rc}{\mathcal{R}}
\newcommand{\View}{\mathsf{View}}
\newcommand{\Adv}{\mathsf{Adv}}
\newcommand{\TV}{\mathrm{TV}}
\newcommand{\call}[1]{\textsf{#1}}

\begin{document}
\lrmaketitle

%%----------------------------------------------------------------------------
\lrscope
The question is whether a learning oracle can answer training and evaluation
requests in unit time without changing any adversary's view. A second question
is where such an oracle becomes useless and where it solves the adversary's
problem outright. We state the problem for any procedure that trains or queries
models, then give a definition, three proven properties and an approximation by
a model bank.

%%----------------------------------------------------------------------------
\lrresult
An oracle that samples exactly from the learning function, costs one unit per
call, takes a record-indexed seed and obeys an access policy is well defined and
lies between useless and trivial. Replacing every training and evaluation step
by an oracle call leaves every adversary's view unchanged in law, running time
excluded (proven, Proposition~1). Naming an honest dataset with its seed decides
membership in one call, and naming only one's own data gains nothing without a
call budget (proven, Propositions~2 and~3). The model's content therefore lies in
the call budget and the access policy, and one policy, neighbouring datasets
with secret seeds, is the adversary of differential privacy. For the
experimenter, who may name everything, the oracle turns a security game into a
sampling procedure with two forms of the leave-one-out model. A bank of models
trained ahead of time answers every call it covers exactly, and errs in
coverage, sampling and hardware nondeterminism.

\lrboundary{This is a definition with proof sketches, and no experiment has
run. We have not checked its novelty against three lines of work, all cited from
memory and not placed. Oracle-efficient learning assumes an empirical risk
minimization oracle (Agarwal et al., ICML 2014). Membership attacks train shadow
models (Carlini et al., IEEE S\&P 2022), and datamodels fit models on many
subsets (Ilyas et al., ICML 2022). A result in the oracle model transfers to a
real procedure only as far as that procedure exposes its node values and no
more. The random-oracle model carries the same caveat (Canetti, Goldreich and
Halevi, STOC 1998, from memory).}

%%----------------------------------------------------------------------------
\lrproduced
\begin{lrfacts}
  \lrfact{Sources}{Halil's statement of the idea in the session of 2026-10-16, and the adversary interface of \lrpath{lib/reports/2026-10-16-halil-adversary-interface.pdf}}
  \lrfact{Artifacts}{this report alone}
  \lrfact{Reproduce}{no code, the source compiles on its own}
\end{lrfacts}

%%----------------------------------------------------------------------------
\lrchanges
The protocol graph of the interface report takes its canonical name, the
dataflow graph, and a semantics in which each learning node is one oracle call.
A question about an adversary becomes a question about its call budget under an
access policy, and an experiment becomes a replay of the dataflow graph against
a model bank. Filling the bank takes as many sequential waves of training as the
graph has learning nodes on its longest path, and each wave trains in parallel.
The next steps are to fix the access policy of each setting we study and to
measure on one small procedure how a membership estimate converges with the
size of the bank.

%%----------------------------------------------------------------------------
\lrmethod
\subsection*{The problem}
Many systems train and use models as steps of a larger procedure. A company
trains a classifier on its users' records and answers prediction queries. A group
of hospitals trains one model across sites. A provider deletes a user's records
and must show that its model forgot them. A data market pays each owner by how
much its data improved a model. Each procedure is a sequence of steps, and each
step trains a model, queries a model, or computes something ordinary, such as an
average, a ciphertext, a noise draw or a release. Training is a learning
function, which maps a starting model, a dataset and a hypothesis class to a
model. Querying is an evaluation map, which maps a model and an input to an
output. We draw the procedure as a directed acyclic graph whose nodes are steps
and whose edges carry datasets and models, and we call it the procedure's
\emph{dataflow graph}, the name the object has in systems work.

An adversary sees some node values, each through a channel: the plain value, a
sum, a ciphertext, a noisy value, or nothing. It sets some inputs and pursues one
of three goals. It distinguishes two hypotheses about the datasets, inverts a
target, or steers the trained model. A proof or a measurement of its success
needs the law of what it sees, and training blocks both. No closed form gives
the law of a trained model. A proof therefore treats training as an arbitrary
function and loses everything specific to learning, or it assumes a structure
that real training lacks, such as a convex loss. Every sample of that law costs a
full training run, and a membership advantage needs many runs per target record.
Experiments shrink as a result. \lrv{MnistCifar} of the \lrv{SweptPapers}
papers of our sweep evaluate on MNIST or CIFAR-10.

The idea is the one cryptography applied to hash functions in the random-oracle
model (Bellare and Rogaway, CCS 1993, from memory). An oracle replaces the
learning function. It answers a training request in one unit of time, with a
model distributed exactly as real training would produce it, and it answers
evaluation requests the same way. A caller that may name any sub-dataset and fix
the randomness obtains, for instance, the model trained on everything except
record $j$ under the randomness of the original run, written $m_{\bar{x}_j}$. A
procedure can then be specified and analysed step by step without the inside of
training, and simulated by looking answers up. The oracle must not change what
the adversary sees when it replaces real training.

Defined loosely, the oracle is trivial or useless. The oracle is trivial when
the adversary may request the model trained on an honest dataset without record
$j$ under the honest randomness. Comparing that model with the released one
decides membership. The oracle is useless when it computes only what the
adversary could compute itself. The best attack of an adversary without resource
limits depends on the law of its view alone, and such an oracle leaves that law
unchanged. The problem is to fix every property that decides between these
two failures: what the oracle computes, whose randomness it uses, what a call
costs, who may name which datasets, which models a caller may open or only
query, and what the oracle refuses. Such an oracle cannot exist, because
training takes time. We approximate it by training the needed models ahead of
time and asynchronously, storing them in a bank keyed by dataset and seed, and
replaying the procedure against the bank.

\subsection*{The setting}
Let $\Z$ be a data space. A \emph{pool} is a finite sequence
$P=(z_1,\dots,z_n)$ of records, and a reference $I\subseteq[n]$ names the dataset
$P_I=\{z_i: i\in I\}$. A dataset may also be given by value, as a finite multiset
over $\Z$. Each record of any dataset carries a key, its index in the pool or a
fresh key assigned when the dataset is submitted. Let $\Mc$ be a model space and
$\HH$ a finite family of named hypothesis classes $\Hc\subseteq\Mc$. A learning
function is a map $\Lf(m,D,\Hc;\rho)\in\Hc$ of a starting model, a dataset, a
class and a seed $\rho\in\Rc$ with law $\mu$. The randomized learning function
of the interface report is the law of $\Lf(m,D,\Hc;\rho)$ under $\rho\sim\mu$.
An evaluation map is $\ev(m,u;\sigma)$ with a seed $\sigma\sim\nu$.

The seed is indexed by record. Write $\rho=(\rho_0,(\rho_k)_{k\in K})$ over the
set $K$ of keys, with $\rho_0\sim\mu_0$ and the $\rho_k$ independent with law
$\mu_1$. The learning function is \emph{local} when it uses $\rho_k$ only in
decisions about record $k$ and $\rho_0$ for everything else, such as the
initialization and noise not tied to a record. Poisson subsampling, where record
$k$ enters step $t$ when coordinate $t$ of $\rho_k$ falls below the rate, is
local. A shuffle that sorts records by a key drawn from $\rho_k$ is local. A
shuffle by one permutation of positions is not, since removing one record moves
every later one.

A procedure is a dataflow graph $G=(V,E)$. Each node has an owner and one of
four operations. An input node holds a dataset or a public object. A learning
node computes $\Lf(\cdot,\cdot,\Hc_v;\rho_v)$ and an evaluation node
$\ev(\cdot,u_v;\sigma_v)$. A general node computes a randomized function, such as
an aggregation, an encryption, a noise draw or a release. An execution evaluates the nodes in
topological order, with independent seeds per node unless the procedure shares
them. An adversary is a coalition with a view, choices and a goal, as in the
interface report.

\subsection*{The oracle}
Definition 1 (learning oracle). Fix $(\Lf,\mu)$, $(\ev,\nu)$, a pool $P$ and the
family $\HH$. A learning oracle $\Or$ keeps a table $T$ from handles to models
and answers three calls.
\begin{itemize}
\item $\call{Train}(a,D,\Hc,\rho)$, with $a$ a model or a handle, $D$ a reference
  or a dataset by value, $\Hc\in\HH$, and $\rho$ a seed or $\bot$. On $\bot$,
  $\Or$ draws $\rho\sim\mu$ and keeps it. $\Or$ stores $\Lf(a,D,\Hc;\rho)$ under a
  fresh handle $h$ and returns $h$.
\item $\call{Eval}(h,u,\sigma)$ returns $\ev(T[h],u;\sigma)$, drawing
  $\sigma\sim\nu$ on $\bot$.
\item $\call{Open}(h)$ returns $T[h]$.
\end{itemize}
It satisfies six properties.
\begin{itemize}
\item Fidelity. Every answer equals the learning function or the evaluation
  map at its seed, and answers with drawn seeds are independent given their
  inputs.
\item Consistency. Equal calls with equal explicit seeds return equal models.
\item Locality. The learning function is local in the record-indexed seed.
\item Unit cost. A call costs one, whatever the dataset size, class or model.
\item Access. Each caller $c$ holds a policy $\Pi_c$ listing the references
  it may name, the seeds it may pass (its own and those $\Pi_c$ makes public),
  the handles it may open, and optionally a sampler of fresh records from a data
  distribution. A handle reaches a caller only as a value it sees, and a caller
  may evaluate every handle it holds. $\Or$ refuses calls outside the policy.
\item Closure. $\Or$ answers no other call. It reveals no seed it drew, no
  state inside training, no preimage of a model, and no running time.
\end{itemize}
A caller that holds a handle it may not open has black-box access, and
$\call{Open}$ gives white-box access. A caller's \call{Eval} calls are counted
apart from its \call{Train} calls, as the interface report counts label queries
against the allowance $Q$.

Definition 2 (learning-oracle model). The procedure $G^{\Or}$ makes each learning
node a \call{Train} call and each evaluation node an \call{Eval} call by the
node's owner, with seeds drawn by $\Or$ unless the procedure fixes them. An
adversary is a coalition with a policy $\Pi$ and a budget $(q_T,q_E)$ of its own
\call{Train} and \call{Eval} calls. The experimenter holds the full policy,
every reference, every seed and every handle.

\subsection*{Three properties}
Proposition 1 (perfect substitution). For every procedure $G$, inputs and
adversary, the view in $G$ with $\Lf$ and $\ev$ computed directly has the same
law as the view in $G^{\Or}$, when the view excludes running time. Proven, by
induction over a topological order of $G$. By fidelity each node value has the same
law given its parents in both executions, and seeds of different nodes are
independent in both. The adversary's own calls have the law of the computation
it would run itself, by fidelity, and closure keeps every other value out of the view. The
exclusion of time is necessary, because unit cost makes a call's time independent of
the dataset size and real training does not.

Proposition 2 (reference access). Let the honest learning node train on
$P_{I\cup\{j\}}$ under $H_1$ and on $P_I$ under $H_0$, with a seed $\rho$ drawn
by $\Or$, and let the adversary see the model.
\begin{itemize}
\item[(a)] If $\Pi$ lets it name $I$ and pass $\rho$, and
  $\Lf(m,P_I,\Hc;\rho)\neq\Lf(m,P_{I\cup\{j\}},\Hc;\rho)$, it wins with advantage
  1 and one call.
\item[(b)] If $\Pi$ lets it name $I$ and $I\cup\{j\}$ with drawn seeds only, its
  advantage is at most $\TV(\View\mid H_0,\View\mid H_1)$ for every budget. When
  the view takes finitely many values, a likelihood-ratio test built from
  $q_T$ samples per hypothesis reaches this bound as $q_T\to\infty$.
\end{itemize}
Proven. Part (a) compares two models. In part (b) the bound is the
data-processing inequality, and the test is optimal by Neyman and Pearson once
its empirical frequencies converge, which they do almost surely. Policy (b) is
the adversary of differential privacy, which knows both neighbouring datasets
and not the seed. The oracle model contains that threat model as one policy, and
a differentially private learning function bounds the advantage of every
adversary under it.

Proposition 3 (own data only). If $\Pi$ names no reference and passes only seeds
the adversary drew, an adversary without a budget gains nothing from $\Or$.
Proven, since every answer is a value of $\Lf$ or $\ev$ on inputs it holds. A
result in the model is therefore an advantage as a function of $(q_T,q_E)$ under
a stated policy. Shadow-model membership attacks are
adversaries of this kind, with a sampler of records from a data distribution
and $q_T$ shadow models.

\subsection*{Two leave-one-out models}
Let $m=\Lf(m_0,P_I,\Hc;\rho)$ with $j\in I$. The coupled leave-one-out model is
$m_{\bar{x}_j}=\Lf(m_0,P_{I\setminus\{j\}},\Hc;\rho)$, under the same seed. The
independent one is $m'_{\bar{x}_j}=\Lf(m_0,P_{I\setminus\{j\}},\Hc;\rho')$, under
a fresh seed $\rho'$. The two have the same law and differ in their joint law
with $m$. Under locality the coupled pair differs only through the
participation of record $j$. It therefore measures what record $j$ did in this
run, the object of influence and memorization measures. The independent pair compares
two laws, which a membership game depends on and differential privacy bounds.
Without locality the coupled form loses its meaning, because removing a record
changes every minibatch.

The experimenter uses the oracle in three ways. It estimates an adversary's
advantage by drawing the hypothesis, running $G^{\Or}$ and the adversary, and
repeating, with a binomial confidence interval. It tests a claimed equality of
laws, such as exact unlearning, which claims that the unlearned model has the law
of $m'_{\bar{x}_j}$. It computes models on many subsets, which data valuation
needs.

\subsection*{Approximation by a model bank}
Definition 3 (model bank). A design $\Delta$ is a finite set of \call{Train}
calls with explicit seeds. The bank $B$ for $\Delta$ stores the model of each call,
computed by running $\Lf$. The empirical oracle $\Ob_B$ answers a call in $\Delta$
by lookup, and \call{Eval} and \call{Open} on stored models. It refuses a call
outside $\Delta$ in a fixed-design experiment, and queues it as a training job in
an asynchronous one.

Every answer of $\Ob_B$ is a sample of $\Lf$ at its seed. Fidelity therefore
holds on $\Delta$ by construction, and unit cost holds once the bank is paid
for. The bank errs in
three ways. Consistency needs deterministic training, and GPU kernels are
nondeterministic by default. Coupled pairs degrade unless the training enforces
deterministic kernels. An adversary whose calls depend on earlier answers may leave
$\Delta$. A fixed bank therefore serves non-adaptive adversaries, and the
asynchronous mode charges adaptivity in waiting time. An estimate from $N$ executions carries
a sampling error of order $N^{-1/2}$, and a true-positive rate at a low
false-positive rate needs far more executions.

A design that includes each record independently with probability one half in
each of $k$ runs gives every record about $k/2$ runs with it and $k/2$ without.
One bank thus serves all $n$ records (the design of Carlini et al., from
memory).
Exact leave-one-out costs $n+1$ runs per seed. A learning node whose starting
model is another learning node's output waits for it. With enough workers one
execution fills in as many waves as its longest path has learning nodes, and
each wave trains in parallel. Federated averaging over $T$ rounds and $N$ clients
takes $T$ waves of $N$ jobs, and a one-shot procedure takes one wave.

%%----------------------------------------------------------------------------
\lrnext
The open choices, each a decision for us.
\begin{itemize}
\item The access policy of each setting we study, from records by value for a
  client to neighbouring references for the adversary of differential privacy.
\item The cost unit, one per call or proportional to the dataset size, which
  prices shadow training on large datasets.
\item The granularity of a learning node, one node per optimizer step, which puts
  the iterates in the graph, or one node per training run, which hides them.
\item Adaptive adversaries against a bank, and how long they may wait.
\item Deterministic training on the cluster, which consistency requires.
\item A public pretrained model as an input node with its own bank entry.
\end{itemize}

\end{document}
